\subsection*{導入：要求が重要な理由}
\addcontentsline{toc}{subsection}{導入：要求が重要な理由}

ソフトウェア要求は、二つの視点から理解する。第一に、要求は「ソフトウェア製品やプロジェクトに対するニーズと制約の表現」である。第二に、要求は「そのニーズと制約を作り、保ち、変化に対応する活動」である。

要求を誤ると、プロジェクトの費用、納期、品質に大きく影響する。ひとつの要求は複数の設計判断を生み、ひとつの設計判断は複数のコード判断を生み、さらに複数のテスト判断につながる。したがって、要求の欠落や誤解が後工程で見つかると、修正範囲が連鎖的に広がる。

現実のプロジェクトで特に多い問題は、\term{不完全性}と\term{曖昧性}である。不完全性とは、必要な要求や詳細が明らかになっていない状態である。曖昧性とは、同じ要求文を複数の意味に解釈できる状態である。どちらも、作ったものが「本来ほしかったもの」とずれる原因になる。

要求文書は、初期開発のためだけではない。保守担当者が後からコードを見ると、コードが「何をしているか」は実行や解析でわかる場合がある。しかし、コードが「何を意図しているか」は要求がなければわかりにくい。欠陥とは、多くの場合、ソフトウェアが意図されたことと実際にしていることの差である。要求は、その意図を長期にわたって伝える役割を持つ。

\MiniBox{重要}{要求活動は、開発の最初に一回だけ行う作業ではない。要求は、プロジェクト開始時に始まり、開発中に洗練され、運用・保守の期間にも参照・変更される。}
